% =====================================================================
%  Proyecto Final — Tema A: Resolución de integrales dobles
%  Análisis de Técnicas Numéricas — CECAR
% ---------------------------------------------------------------------
%  Archivo listo para compilar en Overleaf (compilador: pdfLaTeX).
%  Sube a Overleaf las 7 imágenes de tema_a/graficas/ para que las
%  gráficas se inserten solas. Si no las subes, se muestra un texto
%  indicándolo y el documento igual compila.
% =====================================================================
\documentclass[12pt,a4paper]{article}

% ------------------------------ Paquetes ------------------------------
\usepackage[utf8]{inputenc}
\usepackage[T1]{fontenc}
\usepackage[spanish,es-tabla]{babel}
\usepackage{amsmath,amssymb,amsfonts}
\usepackage[margin=2.5cm]{geometry}
\usepackage{graphicx}
\usepackage{booktabs}
\usepackage{array}
\usepackage{enumitem}
\usepackage{caption}
\usepackage{float}
\usepackage{xcolor}
\usepackage[hidelinks]{hyperref}

% --------------------------- Rutas de imágenes -------------------------
\graphicspath{{./}{graficas/}{../tema_a/graficas/}}

% -------------------------------- Atajos -------------------------------
\newcommand{\dd}{\,\mathrm{d}}

% Figura robusta: si falta la imagen, compila igual mostrando un aviso.
\newcommand{\fig}[4]{%
  \begin{figure}[H]\centering
    \IfFileExists{#2}%
      {\includegraphics[width=#1\linewidth]{#2}}%
      {\fbox{\parbox{0.72\linewidth}{\small\emph{Sube la imagen
        \texttt{#2} a Overleaf (carpeta \texttt{graficas/}) para ver la
        gráfica.}}}}
    \caption{#3}\label{#4}
  \end{figure}}

% ------------------------------ Documento ------------------------------
\title{\textbf{Resolución y análisis comparativo de integrales dobles}\\[4pt]
       \large mediante métodos numéricos\\[6pt]
       \normalsize Proyecto Final — Tema A}
\author{Keiner Elías Beltrán Vergara}
\date{\today}

\begin{document}

% ------------------------------ Portada --------------------------------
\begin{titlepage}
  \centering
  {\scshape\large Corporación Universitaria del Caribe (CECAR)\par}
  {\scshape Facultad de Ciencias Básicas, Ingeniería y Arquitectura\par}
  \vfill
  {\Huge\bfseries Análisis de Técnicas Numéricas\par}
  \vspace{1cm}
  {\LARGE Proyecto Final — Tema A\par}
  \vspace{0.6cm}
  {\large Resolución y comparación de métodos numéricos\par
         para integrales dobles\par}
  \vfill
  \begin{tabular}[t]{l}
    \textbf{Autor:} Keiner Elías Beltrán Vergara\\[4pt]
    \textbf{Docente:} Ing. Carlos Cohen M.\\[4pt]
    \textbf{Asignatura:} Análisis de Técnicas Numéricas\\[4pt]
    \textbf{Fecha:} \today
  \end{tabular}
  \vfill
\end{titlepage}

\tableofcontents
\newpage

% ------------------------------ Resumen --------------------------------
\section*{Resumen}
\addcontentsline{toc}{section}{Resumen}
En este proyecto se resuelven tres integrales dobles de forma
\textbf{analítica} (integración iterada) y \textbf{numérica}, empleando los
métodos vistos en clase: trapecio compuesto, Simpson 1/3, Simpson 3/8,
integración de Romberg y cuadratura de Gauss--Legendre. Para cada método se
calcula el \textbf{error absoluto} y \textbf{relativo} respecto del valor
exacto, se generan gráficas comparativas y de convergencia, y se determina el
método más eficiente. Los resultados se obtienen con cálculos manuales, con
hoja de cálculo (Excel) y con un \textbf{software de propósito general}
desarrollado en Python (consola e interfaz gráfica). Se concluye que
Gauss--Legendre y Romberg alcanzan la solución exacta (error de redondeo de
máquina) para los integrandos polinomiales propuestos, mientras que el
trapecio converge con orden $O(h^2)$.

% --------------------------- Introducción ------------------------------
\section{Introducción}
La integración numérica (o cuadratura) permite aproximar integrales definidas
cuando no existe una primitiva cerrada o cuando el integrando solo se conoce
en un conjunto discreto de puntos. En este trabajo se estudia el caso de
\textbf{integrales dobles} sobre regiones rectangulares, que se reducen a dos
integrales simples mediante \textbf{integración iterada}:
\[
  \iint_R f(x,y)\dd x\dd y
  = \int_{y_a}^{y_b}\!\left[\int_{x_a}^{x_b} f(x,y)\dd x\right]\dd y .
\]
Se implementan cinco familias de métodos (con variantes) y se compara su
eficiencia a partir del error absoluto cometido frente al valor analítico
exacto.
% ------------------------------ Objetivos ------------------------------
\section{Objetivos}

\subsection{Objetivo general}
Resolver tres integrales dobles de forma analítica y numérica, empleando todos
los métodos vistos en clase, y comparar la eficiencia de cada método con base
en su error absoluto.

\subsection{Objetivos específicos}
\begin{itemize}[leftmargin=1.4em]
  \item Calcular analíticamente (manualmente) el valor exacto de cada integral
        mediante integración iterada.
  \item Implementar los métodos de trapecio, Simpson 1/3, Simpson 3/8,
        Romberg y Gauss--Legendre en un software de propósito general.
  \item Reproducir los cálculos en una hoja de cálculo (Excel) con fórmulas
        dinámicas.
  \item Determinar el error absoluto y relativo de cada método y clasificarlos
        según su eficiencia.
  \item Analizar la convergencia de los métodos (orden de exactitud y
        coeficiente de determinación $R^2$ del ajuste log--log).
\end{itemize}

% ---------------------------- Marco teórico -----------------------------
\section{Marco teórico}
Sea $I=\int_a^b f(x)\dd x$. En todos los métodos se particiona $[a,b]$ en $n$
subintervalos de longitud $h=(b-a)/n$ y se aproxima la integral con una suma
ponderada de valores de $f$.

\subsection{Regla del trapecio compuesta}
Aproxima $f$ por tramos rectos (polinomios de grado 1):
\[
  T_n = \frac{h}{2}\Big[f(x_0)+2\sum_{i=1}^{n-1}f(x_i)+f(x_n)\Big],
  \qquad \text{error } O(h^2).
\]
Al duplicar $n$, el error se reduce aproximadamente a la cuarta parte.

\subsection{Regla de Simpson 1/3 compuesta}
Aproxima $f$ por parábolas (polinomios de grado 2) en pares de subintervalos
($n$ par):
\[
  S_n = \frac{h}{3}\Big[f(x_0)+4\!\!\sum_{i\ \text{impar}}\!\!f(x_i)
        +2\!\!\sum_{i\ \text{par}}\!\!f(x_i)+f(x_n)\Big],
  \qquad \text{error } O(h^4).
\]
Integra exactamente polinomios de grado $\le 3$.

\subsection{Regla de Simpson 3/8 compuesta}
Aproxima $f$ por cúbicas en grupos de tres subintervalos ($n$ múltiplo de 3):
\[
  S_{3/8} = \frac{3h}{8}\big[f(x_0)+f(x_n)
        +3\!\!\sum_{i\neq0,n,\ i\not\equiv 0 (3)}\!\!f(x_i)
        +2\!\!\sum_{i\equiv 0 (3)}\!\!f(x_i)\big],
  \qquad \text{error } O(h^4).
\]

\subsection{Integración de Romberg}
Combina la regla del trapecio con la \textbf{extrapolación de Richardson} para
eliminar términos de error sucesivos, obteniendo aproximaciones de orden
creciente $O(h^{2m})$ con $m$ niveles. Su elemento final $R_{m,m}$ es la mejor
aproximación de la tabla.

\subsection{Cuadratura de Gauss--Legendre}
Elige de forma óptima los nodos $x_i$ (raíces del polinomio de Legendre) y los
pesos $w_i$ para maximizar la exactitud:
\[
  G_n = \frac{b-a}{2}\sum_{i=1}^{n} w_i\, f\!\left(\frac{a+b}{2}
        +\frac{b-a}{2}x_i\right).
\]
Con $n$ puntos integra \textbf{exactamente} cualquier polinomio de grado
$\le 2n-1$. Por ello es el método más eficiente para integrandos polinomiales.

\subsection{Integrales dobles por integración iterada}
Para $R=[x_a,x_b]\times[y_a,y_b]$ se aplica un método 1D en $x$ (integral
interior) y el mismo método en $y$ (integral exterior):
\[
  \iint_R f(x,y)\dd x\dd y
  = \int_{y_a}^{y_b} g(y)\dd y,\qquad
  g(y)=\int_{x_a}^{x_b} f(x,y)\dd x .
\]
% ------------------------ Planteamiento del problema --------------------
\section{Planteamiento del problema}
Se pide resolver las siguientes tres integrales dobles de forma analítica y
numérica:

\begin{center}
\begin{tabular}{c c c c}
\toprule
 & Integral & Región $R$ & Valor exacto \\
\midrule
1 & $\displaystyle\iint (1-x^2-y^2)\dd x\dd y$ & $[-1,1]\times[-1,1]$ & $4/3$ \\[4pt]
2 & $\displaystyle\iint x^2 y^2\dd x\dd y$ & $[1,2]\times[1,3]$ & $182/9$ \\[4pt]
3 & $\displaystyle\iint (x^2+y^2)\dd x\dd y$ & $[-1,1]\times[-1,1]$ & $8/3$ \\
\bottomrule
\end{tabular}
\end{center}

\noindent\textit{Nota:} el enunciado también menciona un valor de referencia
``exacto'' de $2.230985$ para un ejercicio de Gauss--Legendre; dicho valor no
coincide con las tres integrales anteriores y debe verificarse visualmente en
el PDF original.

% --------------------------- Solución analítica -------------------------
\section{Solución analítica (cálculo manual)}

\subsection{Integral 1: $\iint (1-x^2-y^2)\dd x\dd y$ sobre $[-1,1]^2$}
Integrando primero respecto de $x$:
\[
  \int_{-1}^{1}(1-x^2-y^2)\dd x
  =\Big[x-\tfrac{x^3}{3}-y^2x\Big]_{-1}^{1}
  =\Big(\tfrac{2}{3}-y^2\Big)-\Big(-\tfrac{2}{3}+y^2\Big)
  =\frac{4}{3}-2y^2 .
\]
Integrando ahora respecto de $y$:
\[
  \int_{-1}^{1}\Big(\frac{4}{3}-2y^2\Big)\dd y
  =\Big[\frac{4y}{3}-\frac{2y^3}{3}\Big]_{-1}^{1}
  =\frac{2}{3}-\Big(-\frac{2}{3}\Big)
  =\frac{4}{3}\approx 1.3333333333 .
\]

\subsection{Integral 2: $\iint x^2 y^2\dd x\dd y$ sobre $[1,2]\times[1,3]$}
Por ser separable (producto de funciones de una variable):
\[
  \int_{1}^{3}\!\int_{1}^{2}x^2y^2\dd x\dd y
  =\Big(\int_{1}^{2}x^2\dd x\Big)\Big(\int_{1}^{3}y^2\dd y\Big)
  =\Big[\frac{x^3}{3}\Big]_{1}^{2}\Big[\frac{y^3}{3}\Big]_{1}^{3}
  =\frac{7}{3}\cdot\frac{26}{3}
  =\frac{182}{9}\approx 20.2222222222 .
\]

\subsection{Integral 3: $\iint (x^2+y^2)\dd x\dd y$ sobre $[-1,1]^2$}
Integrando primero respecto de $x$:
\[
  \int_{-1}^{1}(x^2+y^2)\dd x
  =\Big[\frac{x^3}{3}+y^2x\Big]_{-1}^{1}
  =\Big(\frac{1}{3}+y^2\Big)-\Big(-\frac{1}{3}-y^2\Big)
  =\frac{2}{3}+2y^2 .
\]
Integrando respecto de $y$:
\[
  \int_{-1}^{1}\Big(\frac{2}{3}+2y^2\Big)\dd y
  =\Big[\frac{2y}{3}+\frac{2y^3}{3}\Big]_{-1}^{1}
  =\frac{4}{3}-\Big(-\frac{4}{3}\Big)
  =\frac{8}{3}\approx 2.6666666667 .
\]
% --------------------------- Solución numérica --------------------------
\section{Solución numérica}
Los cálculos se realizaron con el software de propósito general desarrollado en
Python (módulo \texttt{metodos.py}), aplicando cada método 1D a la integral
interior (en $x$) y a la exterior (en $y$). Los parámetros utilizados son:
$n=100$ subintervalos para trapecio y Simpson 1/3, $n=99$ para Simpson 3/8
(múltiplo de 3), $m=8$ niveles para Romberg y $n=2,3,5$ puntos para
Gauss--Legendre. Los resultados se muestran a continuación.

\subsection{Integral 1: $I_1=4/3\approx 1.3333333333$}
\begin{center}
\begin{tabular}{l c c c}
\toprule
Método & Resultado & Error absoluto & Error relativo \\
\midrule
Trapecio ($n=100$)       & $1.3328000000$ & $5.333\times10^{-4}$ & $4.000\times10^{-4}$ \\
Simpson 1/3 ($n=100$)    & $1.3333333333$ & $2.220\times10^{-16}$ & $1.665\times10^{-16}$ \\
Simpson 3/8 ($n=99$)     & $1.3333333333$ & $4.441\times10^{-16}$ & $3.331\times10^{-16}$ \\
Romberg ($m=8$)          & $1.3333333333$ & $1.110\times10^{-15}$ & $8.327\times10^{-16}$ \\
Gauss--Legendre (2 ptos) & $1.3333333333$ & $4.441\times10^{-16}$ & $3.331\times10^{-16}$ \\
Gauss--Legendre (3 ptos) & $1.3333333333$ & $0$ & $0$ \\
Gauss--Legendre (5 ptos) & $1.3333333333$ & $1.110\times10^{-15}$ & $8.327\times10^{-16}$ \\
\bottomrule
\end{tabular}
\end{center}
\textbf{Método más eficiente:} Gauss--Legendre (3 puntos), error $=0$.

\subsection{Integral 2: $I_2=182/9\approx 20.2222222222$}
\begin{center}
\begin{tabular}{l c c c}
\toprule
Método & Resultado & Error absoluto & Error relativo \\
\midrule
Trapecio ($n=100$)       & $20.2226777800$ & $4.556\times10^{-4}$ & $2.253\times10^{-5}$ \\
Simpson 1/3 ($n=100$)    & $20.2222222222$ & $3.553\times10^{-15}$ & $1.757\times10^{-16}$ \\
Simpson 3/8 ($n=99$)     & $20.2222222222$ & $1.066\times10^{-14}$ & $5.271\times10^{-16}$ \\
Romberg ($m=8$)          & $20.2222222222$ & $0$ & $0$ \\
Gauss--Legendre (2 ptos) & $20.2222222222$ & $0$ & $0$ \\
Gauss--Legendre (3 ptos) & $20.2222222222$ & $3.553\times10^{-15}$ & $1.757\times10^{-16}$ \\
Gauss--Legendre (5 ptos) & $20.2222222222$ & $0$ & $0$ \\
\bottomrule
\end{tabular}
\end{center}
\textbf{Método más eficiente:} Romberg ($m=8$), error $=0$ (junto con
Gauss--Legendre de 2 y 5 puntos, que también dan error $0$).

\subsection{Integral 3: $I_3=8/3\approx 2.6666666667$}
\begin{center}
\begin{tabular}{l c c c}
\toprule
Método & Resultado & Error absoluto & Error relativo \\
\midrule
Trapecio ($n=100$)       & $2.6672000000$ & $5.333\times10^{-4}$ & $2.000\times10^{-4}$ \\
Simpson 1/3 ($n=100$)    & $2.6666666667$ & $8.882\times10^{-16}$ & $3.331\times10^{-16}$ \\
Simpson 3/8 ($n=99$)     & $2.6666666667$ & $8.882\times10^{-16}$ & $3.331\times10^{-16}$ \\
Romberg ($m=8$)          & $2.6666666667$ & $0$ & $0$ \\
Gauss--Legendre (2 ptos) & $2.6666666667$ & $0$ & $0$ \\
Gauss--Legendre (3 ptos) & $2.6666666667$ & $1.332\times10^{-15}$ & $4.996\times10^{-16}$ \\
Gauss--Legendre (5 ptos) & $2.6666666667$ & $4.441\times10^{-16}$ & $1.665\times10^{-16}$ \\
\bottomrule
\end{tabular}
\end{center}
\textbf{Método más eficiente:} Romberg ($m=8$), error $=0$ (junto con
Gauss--Legendre de 2 puntos, que también da error $0$).
% ---------------------- Comparación de eficiencia ----------------------
\section{Comparación de eficiencia}
La eficiencia se evalúa con el \textbf{error absoluto} $|I_{\text{num}}-I_{\text{exacto}}|$
frente al número de evaluaciones del integrando. De las tablas anteriores se
observa:

\begin{itemize}[leftmargin=1.4em]
  \item \textbf{Trapecio} es el método menos preciso: su error es del orden de
        $5\times10^{-4}$ con $n=100$ y converge con orden $O(h^2)$.
  \item \textbf{Simpson 1/3 y 3/8} alcanzan la precisión de máquina
        ($\sim10^{-15}$) porque los integrandos son polinomios de grado bajo y
        ambos métodos son exactos para grado $\le 3$.
  \item \textbf{Romberg} y \textbf{Gauss--Legendre} llegan al valor exacto con
        error $0$ o del orden de la unidad de redondeo ($\sim10^{-16}$),
        confirmando que son los métodos más eficientes para estos integrandos.
\end{itemize}

\subsection{Convergencia del trapecio y orden de exactitud}
La siguiente tabla muestra el error absoluto del trapecio al aumentar el número
de subintervalos $n$ en cada dirección:

\begin{center}
\begin{tabular}{c c c c}
\toprule
$n$ & Error ($I_1$) & Error ($I_2$) & Error ($I_3$) \\
\midrule
4   & $3.333\times10^{-1}$ & $2.856\times10^{-1}$ & $3.333\times10^{-1}$ \\
8   & $8.333\times10^{-2}$ & $7.123\times10^{-2}$ & $8.333\times10^{-2}$ \\
16  & $2.083\times10^{-2}$ & $1.780\times10^{-2}$ & $2.083\times10^{-2}$ \\
32  & $5.208\times10^{-3}$ & $4.449\times10^{-3}$ & $5.208\times10^{-3}$ \\
64  & $1.302\times10^{-3}$ & $1.112\times10^{-3}$ & $1.302\times10^{-3}$ \\
128 & $3.255\times10^{-4}$ & $2.780\times10^{-4}$ & $3.255\times10^{-4}$ \\
\bottomrule
\end{tabular}
\end{center}

Al duplicar $n$, el error se reduce a la cuarta parte (factor $\approx 4$), lo
que confirma un orden de convergencia $p=2$. Ajustando $\log_{10}E$ contra
$\log_{10}n$ mediante mínimos cuadrados se obtiene:
\[
  \text{pendiente } = -p = -2.001,\qquad R^2 = 1.0000 .
\]
El coeficiente de determinación $R^2\approx 1$ indica que los datos siguen
perfectamente la ley $E\propto n^{-2}$.

Para Simpson y Gauss--Legendre el error ya está en el nivel de la precisión de
máquina ($\sim10^{-15}$) desde $n$ pequeño, por lo que no existe una tendencia
polinomial que ajustar (el método ya es ``exacto'' a efectos prácticos).

% ------------------------------- Gráficas -------------------------------
\section{Gráficas generadas}
Las gráficas fueron generadas automáticamente por el software
(\texttt{graficas.py}) y guardadas en formato PNG.

\fig{0.85}{barras_integral_1.png}{Error absoluto por método para la integral 1
($1-x^2-y^2$).}{fig:barras1}
\fig{0.85}{barras_integral_2.png}{Error absoluto por método para la integral 2
($x^2y^2$).}{fig:barras2}
\fig{0.85}{barras_integral_3.png}{Error absoluto por método para la integral 3
($x^2+y^2$).}{fig:barras3}
\fig{0.85}{convergencia_integral_1.png}{Convergencia (error vs.\ $n$) para la
integral 1.}{fig:conv1}
\fig{0.85}{convergencia_integral_2.png}{Convergencia (error vs.\ $n$) para la
integral 2.}{fig:conv2}
\fig{0.85}{convergencia_integral_3.png}{Convergencia (error vs.\ $n$) para la
integral 3.}{fig:conv3}
\fig{0.95}{comparacion_global.png}{Comparación global de errores por método e
integral.}{fig:global}

% -------------------------------- Excel ---------------------------------
\section{Cálculos en Excel}
Se elaboró una hoja de cálculo (\texttt{Tema\_A\_Excel.xlsx}) con cuatro hojas
(\emph{Resumen}, \emph{Metodos\_1D}, \emph{Gauss\_Romberg} y
\emph{Metodología}) que reproduce los cálculos anteriores mediante
\textbf{fórmulas en vivo}, de modo que al abrirla los valores se recalculan y
pueden verificarse. Las fórmulas implementan las mismas sumas ponderadas del
marco teórico (trapecio, Simpson y Gauss--Legendre con nodos y pesos tabulados).

% ----------------------------- Conclusiones -----------------------------
\section{Conclusiones}
\begin{itemize}[leftmargin=1.4em]
  \item Todas las integrales pudieron resolverse analíticamente por integración
        iterada, obteniendo $I_1=4/3$, $I_2=182/9$ e $I_3=8/3$.
  \item Los métodos de orden superior (Simpson, Romberg y Gauss--Legendre)
        reproducen el valor exacto con errores del orden de la precisión de
        máquina ($10^{-15}$--$10^{-16}$), por lo que son los más eficientes.
  \item El trapecio converge con orden $O(h^2)$, confirmado por el ajuste
        log--log ($p\approx2$, $R^2=1$), siendo el método menos eficiente.
  \item Gauss--Legendre con $n$ puntos integra exactamente polinomios de grado
        $\le 2n-1$: por eso con 2--3 puntos basta para estos integrandos.
\end{itemize}

% --------------------------- Recomendaciones ----------------------------
\section{Recomendaciones}
\begin{itemize}[leftmargin=1.4em]
  \item Para integrandos suaves y polinomiales, usar Gauss--Legendre o Romberg,
        que requieren pocas evaluaciones y alcanzan alta exactitud.
  \item Evitar el trapecio cuando se necesite alta precisión; en su lugar,
        emplear Simpson (compuesta) o extrapolación de Romberg.
  \item Verificar visualmente los límites de integración en el enunciado
        original y, de ser necesario, actualizar el valor exacto en el código.
\end{itemize}

% ---------------------------- Referencias -------------------------------
\begin{thebibliography}{9}
\bibitem{burden} R.~L. Burden y J.~D. Faires, \emph{Análisis Numérico},
9.ª ed., Cengage Learning, 2011.
\bibitem{chapra} S.~C. Chapra y R.~P. Canale, \emph{Métodos Numéricos para
Ingenieros}, 7.ª ed., McGraw-Hill, 2015.
\bibitem{abramowitz} M.~Abramowitz e I.~A. Stegun, \emph{Handbook of
Mathematical Functions}, Dover, 1972 (nodos y pesos de Gauss--Legendre).
\end{thebibliography}

\end{document}




